\section{Limitations}
\label{sec:limitations}

Every observed difference could partly be an artefact of the following
factors. They are listed so readers can weigh the results accordingly.

\paragraph{Machine.}
One laptop-class host (i7-10750H, 8 GB RAM), under WSL2 virtualisation, with
an NTFS-backed filesystem. CPU frequency scaling, thermal throttling and the
hypervisor layer affect all languages, but not necessarily equally. Every
timed harness phase is pinned to a single logical core with \texttt{taskset}
(\cref{sec:host}), which bounds the CPU budget but is no performance governor
and no bare metal: neighbouring work still shares that core.

\paragraph{Runtime.}
The Java harness and JMH both run on OpenJDK 21 (the harness through
\texttt{mvn exec:java}, whose own JVM overhead is outside the timed regions
but shares the machine; JMH in a shaded jar). Python runs on 3.14, Go 1.27,
Rust 1.99: these are current
versions, not long-term-stability comparisons. Only Java receives warm-up
(\cref{sec:warmup}); the other languages are measured from cold start, which
is the honest default for them but is itself a protocol choice.

\paragraph{Cache state.}
The OS page cache is never dropped between phases, so a file read twice in
one phase may be served from RAM. All languages in a phase share the same
state, but phases are not independent of the ones before them. MongoDB's
internal cache (WiredTiger) is likewise warm after the first build.

\paragraph{Workload.}
The corpus is \fixed{100} books (\fixed{61.4} MB); the scalability
builds reach 10{,}000 ``books'' by repeating that corpus, so index size
grows but term diversity does not. The datalake experiments stay at the same \fixed{100} books: they compare
directory structures on identical work, and growing them would mean
downloading thousands more books (breaking the frozen artifact) or
duplicating files (measuring nothing new) --- no datalake scaling curve was
measured. Queries are \fixed{30} fixed lines; they are
representative of the assignment but not sampled from a query log.
Results may reorder on corpora with different term distributions.

\paragraph{Timing and statistics.}
Elapsed times are wall-clock on a shared host where only the CPU budget is
pinned (\cref{sec:host}): neighbouring
processes (including \texttt{mongod} and the sampling threads) inject
jitter, visible as non-zero \texttt{stdev\_seconds}. One tournament run
produces the committed snapshot; there is no second run to quantify
run-to-run variance. No outliers are trimmed and no phase is repeated for a
better result.

\paragraph{Measurement boundaries.}
Client RSS excludes the MongoDB server by construction
(\cref{sec:index-memory}); its storage is measured via \texttt{collStats},
its memory is not. The download probe is ten sequential requests per round
over one pooled connection, repeated in five interleaved rounds, so its
rate reflects latency more than bandwidth. CPU percentage is process CPU
over wall clock; the pinned harness phases cannot exceed 100.

\paragraph{Unmeasured behaviours.}
Not exercised at all: concurrent readers and writers, indexes far beyond the
measured batches, MongoDB replication/sharding, crash recovery of the
database itself, and the Stage 2 serving path. Conclusions in
\cref{sec:conclusions} do not extend to them.
